\section{Exponential Functions and Logarithms}

Exponential functions and logarithms play a central role in mathematics and physics. They are especially important when describing processes that grow or decay over time, such as sound intensity, population growth, or radioactive decay.

\subsection{Exponential Functions}

Exponential functions describe processes in which a quantity changes by a constant \textbf{percentage} rather than a constant amount. This is a key difference from linear functions. This means: the bigger the quantity already is, the faster it grows (or decays).

A linear function changes by adding the same amount each step:
\begin{equation}
    y = mx + b
\end{equation}

An exponential function changes by multiplying with the same factor each step:
\begin{equation}
    f(x) = a \cdot b^x
\end{equation}

\begin{itemize}
    \item $a$: initial value (starting amount)
    \item $b$: growth factor (multiplication factor per step)
\end{itemize}

\paragraph{Growth Factor and Percent Change}

The growth factor $b$ describes how much a quantity changes per step:
\begin{itemize}
    \item $b > 1$: exponential growth ($f(x)$ increases with increasing $x$ )
    \item $0 < b < 1$: exponential decay ($f(x)$ decreases with increasing $x$ )
\end{itemize}

\paragraph{Graphs of Exponential Functions}
%Plots von E-Funktionen einfügen%
When considering the graph of an exponential function one observes the following properties:

\begin{itemize}
    \item The graph is a smooth curve (not a straight line as e.g. a linear function)
    \item It is always above the x-axis (for $a > 0$)
    \item It has a \textbf{horizontal asymptote}
\end{itemize}

\paragraph{Horizontal Asymptote:} 
A horizontal asymptote is a line that the graph approaches but never reaches. Considering an exponential growth, the graph approches $y=0$ for $x \rightarrow -\infty$. So, with dexreasing $x$, the graph gets closer and closer to zero but never reaches the $x-$axis.



\paragraph{Long-Term Behavior}
\begin{itemize}
    \item For exponential growth ($b>1$): $f(x) \to \infty$ as $x \to \infty$, $f(x) \to 0$ as $x \to - \infty$, 
    \item For exponential  decay ($0<b<1$): $f(x) \to 0$ as $x \to \infty$ 
\end{itemize}

\paragraph{Euler's number $e$}
Euler's number, denoted as $e$, is one of the non-terminating constants we often use in mathematics and science. Another non-terminating constant you might already know is $\pi$ which is often used in geometry.
Euler’s number approximately equals 2.71828. You also obtain $e$ when calculating the following term for very large numbers $n$: 
\begin{align}
    e \approx  \left( 1+ \frac{1}{n} \right)^n
\end{align}
With increasing $n$ the term approches $e$ as you can see in the following table: 
\begin{table}[h]
\centering
\begin{tabular}{|c|c|}
\hline
$n$ & $\left(1 + \frac{1}{n}\right)^n$ \\
\hline
1      & 2{.}000000 \\
2      & 2{.}250000 \\
5      & 2{.}488320 \\
10     & 2{.}593742 \\
100    & 2{.}704814 \\
1000   & 2{.}716924 \\
10000  & 2{.}718146 \\
100000 & 2{.}718268 \\
\hline
$\infty$ & $e \approx 2{.}718282\ldots$ \\
\hline
\end{tabular}
\caption{approximation of Euler's number $e$ as $n$ increases}
\label{tab:euler}
\end{table}
Euler’s number helps mathematicians, physicists and other scientists explain exponential growth in various phenomena, from population to radioactive decay. So, euler's number $e$ often becomes the growth factor when describing natural phenomena. The following function is thus very important in math and physics: 
\begin{align}
    f(x) = a \cdot e^x
\end{align}




\paragraph{Transformations of Exponential Functions}

Similar to sine and cosine functions, exponential functions can also be changed using different parameters. These parameters allow us to adjust the shape and position of the graph.

The general form is:
\begin{equation}
    f(x) = a \cdot b^{k(x-c)} + d
\end{equation}

Each parameter has a clear effect on the graph:

\begin{itemize}
    \item \textbf{Parameter $a$ (vertical scaling):}  
    This parameter stretches or compresses the graph vertically.
    \begin{itemize}
        \item $a > 1$: the graph becomes steeper (grows faster visually)
        \item $0 < a < 1$: the graph becomes flatter
        \item $a < 0$: the graph is flipped upside down (reflection at the x-axis)
    \end{itemize}

    \item \textbf{Parameter $k$ (growth speed):}  
    This parameter controls how fast the function grows or decays.
    \begin{itemize}
        \item larger $k$: faster growth or decay
        \item smaller $k$: slower change
    \end{itemize}
    You can think of $k$ as controlling how ``quickly'' the curve rises or falls.

    \item \textbf{Parameter $c$ (horizontal shift):}  
    This moves the graph left or right.
    \begin{itemize}
        \item shift to the right if $c > 0$
        \item shift to the left if $c < 0$
    \end{itemize}
    This means the function starts growing earlier or later.

    \item \textbf{Parameter $d$ (vertical shift):}  
    This moves the entire graph up or down.
    \begin{itemize}
        \item $d > 0$: graph moves up
        \item $d < 0$: graph moves down
    \end{itemize}
\end{itemize}


By changing these parameters, we can adapt exponential functions to model many real-world processes, such as population growth or decay processes in physics.


\subsection{Logarithms}

The logarithm is the inverse operation of an exponential function. While an exponential function answers the question

\begin{center}
    \textit{``What is the result if we raise a number to a certain power?''}
\end{center}

the logarithm answers the reverse question:

\begin{center}
    \textit{``To which power must we raise a number to obtain a given result?''}
\end{center}

Mathematically, this relationship is written as:
\begin{equation}
    b^x = y \quad \Longleftrightarrow \quad \log_b(y) = x
\end{equation}

Here:
\begin{itemize}
    \item $b$ is the base
    \item $x$ is the exponent
    \item $y$ is the result
\end{itemize}

\paragraph{Example:}
\begin{equation}
    2^3 = 8 \quad \Longleftrightarrow \quad \log_2(8) = 3
\end{equation}

\paragraph{Intuitive Explanation:}
\begin{itemize}
    \item Exponential function: you start with a number and calculate the result
    \item Logarithm: you are given the result and want to find the exponent
\end{itemize}

\paragraph{Important Notes:}
\begin{itemize}
    \item Logarithms are only defined for positive numbers ($y > 0$)
    \item The base must be positive and not equal to 1
\end{itemize}

\paragraph{Summary:}
The logarithm is a tool that helps us ``undo'' exponential growth. It is especially useful when we want to solve equations where the unknown is in the exponent.


\subsubsection{Natural Logarithm}

The logarithm with base $e$ is called the natural logarithm and denoted by $\ln$:
\begin{equation}
     log_e (y) = \ln(y)
\end{equation}


So the natural logarithm answers the question:
\begin{center}
    \textit{``To which power must we raise $e$ to get a certain number?''}
\end{center}


\paragraph{Inverse Relationship}

Exponential and logarithmic functions undo each other:
\begin{equation}
    e^{\ln(x)} = x
\end{equation}
\begin{equation}
    \ln(e^x) = x
\end{equation}

%Plots von ln und e-Fkt einfügen%

From the graphs, we can also see that the natural logarithm function and the exponential function cancel each other out. This means they are \textbf{inverse functions}.

Geometrically, this relationship can be observed by the fact that their graphs are mirror images of each other along the line
\begin{equation}
    y = x
\end{equation}
which is also called the angle bisector.

\paragraph{Intuitive Explanation:}
\begin{itemize}
    \item The exponential function $e^x$ takes an input $x$ and produces a value
    \item The logarithm $\ln(x)$ takes this value and returns the original input $x$
\end{itemize}

So applying one function after the other brings you back to where you started:
\begin{equation}
    \ln(e^x) = x \quad \text{and} \quad e^{\ln(x)} = x
\end{equation}

\subsection{Summary}

\begin{itemize}
    \item Exponential functions describe growth and decay processes
    \item The number $e$ is fundamental for natural growth
    \item Logarithms are the inverse of exponential functions
    \item Both concepts are essential for modeling real-world phenomena
\end{itemize}
